
\section*{Chapitre 9 - Distance et compas}
\addcontentsline{toc}{section}{Chapitre 9 - Distance et compas}

\subsection*{I. Cercles}
\addcontentsline{toc}{subsection}{I. Cercles}


\begin{Définitions*}[c]

\begin{minipage}[c]{0.65\textwidth}
\vspace{0.2cm}
Un \textbf{cercle} est l'ensemble de tous les points situés à une \textbf{même distance} d'un même point appelé \textbf{le centre}. Cette distance est appelée \textbf{le rayon} du cercle.
\begin{itemize}[label=\textcolor{Couleur6}{$\bullet$}]
\item Un \textbf{rayon} est un segment reliant le centre à un point du cercle.
\item Un \textbf{diamètre} est un segment reliant deux points du cercle et passant par le centre.
\item Une \textbf{corde} est un segment reliant deux points distincts du cercle.
\end{itemize}

\end{minipage} \hfill \begin{minipage}[c]{0.3\textwidth}
\begin{center}
\begin{tikzpicture}[scale=1]
    % Cercle
    \draw[very thick, Couleur8] (0,0) circle (2cm);
    

    
    % Points sur le cercle
    \coordinate (A) at (135:2);
    \coordinate (B) at (-45:2);
    \coordinate (C) at (30:2);
    \coordinate (D) at (-120:2);
    \coordinate (O) at (0:0);
    

    
    % Diamètre
    \draw[Couleur11, very thick] (A) -- (B);
    \node[Couleur11, rotate=-45, above] at (-0.7,0.7) {\textbf{Diamètre}};
    
    % Rayon
    \draw[Couleur7, very thick] (0,0) -- (C);
    \node[Couleur7, above right] at (1.8,1) {\textbf{Rayon}};
    
    % Corde
    \draw[Couleur12, very thick] (B) -- (D);
    \node[Couleur12, below] at (0,-2) {\textbf{Corde}};
    
    % Croix sur les points
    \foreach \point in {A,B,C,D,O} {
        \draw[black, thick] (\point) +(-0.1,0.1) -- +(0.1,-0.1);
        \draw[black, thick] (\point) +(-0.1,-0.1) -- +(0.1,0.1);
    }
    
        % Centre

    \node[below] at (0,0) {Centre};
\end{tikzpicture}
\end{center}

\end{minipage}
\end{Définitions*}

\begin{Remarque}
Les termes « rayon » et « diamètre » désignent à la fois un segment et une longueur.
\end{Remarque}

\begin{Propriétés*}
\leavevmode\vspace{-\baselineskip}
\begin{itemize}[label=\textcolor{Couleur1}{$\bullet$}]
\item La longueur d'un diamètre correspond à deux fois celle du rayon. Elle est toujours supérieure ou égale à la longueur d'une corde.
\item La longueur L d'un cercle de diamètre D est \textbf{proportionnelle} à son diamètre. Elle est égale à L = $\pi$ × D ou encore L = 2 × $\pi$ × r où r représente le rayon.
\item Le nombre $\pi$ (qui se lit « pi ») n'est pas un nombre décimal et ne peut pas s'écrire sous la forme d'une fraction, il possède une infinité de décimales. On a $\pi$ $\simeq$ 3,14.
\end{itemize}
\end{Propriétés*}



\begin{Exemple}
La longueur d'un cercle de rayon 7 cm vaut : L = 2 × $\pi$ × 7 cm = 14$\pi\simeq$ 44 cm.
\end{Exemple}


\subsection*{II. Disques}
\addcontentsline{toc}{subsection}{II. Disques}

\begin{Définition}
Le \textbf{disque} de centre O et de rayon r est l'ensemble de tous les points situés à une distance de O \textbf{inférieure ou égale} à r. Il s'agit de la \textbf{surface} délimitée par le cercle de même centre et de même rayon.
\end{Définition}

\begin{Exemple}
On a tracé le cercle et le disque de centre O et de rayon 4 cm. \\
\begin{minipage}[c]{0.4\textwidth}
\begin{enumerate}
\item Compléter par $\in$ ou $\notin$ :
\begin{itemize}
\item Le point C $\in$ au disque.
\item Le point C $\notin$ au cercle.
\item Le point H $\in$ au disque et au cercle.
\item Le point G $\notin$ au disque.
\end{itemize}
\end{enumerate}
\end{minipage} \hfill \begin{minipage}[c]{0.6\textwidth}
\begin{center}
\begin{tikzpicture}[scale=0.8]
    % Disque rempli
    \fill[Couleur8!10] (0,0) circle (2cm);
    
    % Cercle (contour)
    \draw[very thick, Couleur8] (0,0) circle (2cm);
    
    % Centre O
    \coordinate (O) at (0:0);
    \node[above right] at (O) {O};
    
    % Point G sur le cercle
    \coordinate (G) at (45:2.2);
    \node[above right] at (G) {G};
    
    % Point H sur le cercle
    \coordinate (H) at (0:2);
    \node[right] at (H) {H};
    
    % Point C sur le cercle (en bas)
    \coordinate (C) at (-90:1.5);
    \node[above left] at (C) {C};
    

    
    % Croix sur les points
        \draw[black, thick] (O) +(-0.1,0.1) -- +(0.1,-0.1);
            \draw[black, thick] (O) +(-0.1,-0.1) -- +(0.1,0.1);
    \draw[black, thick] (G) +(-0.1,0.1) -- +(0.1,-0.1);
    \draw[black, thick] (G) +(-0.1,-0.1) -- +(0.1,0.1);
    \draw[black, thick] (H) +(-0.1,0.1) -- +(0.1,-0.1);
    \draw[black, thick] (H) +(-0.1,-0.1) -- +(0.1,0.1);
    \draw[black, thick] (C) +(-0.1,0.1) -- +(0.1,-0.1);
    \draw[black, thick] (C) +(-0.1,-0.1) -- +(0.1,0.1);
    
  
\end{tikzpicture}
\end{center}

\end{minipage} 
\begin{enumerate}[start=2]
\item Compléter par $<$ ou $>$ ou $=$ :
\vspace{-0.5cm}
\begin{multicols}{3}
\begin{itemize}
\item OC $<$ 4cm
\item OH $=$ 4cm
\item OG $>$ 4cm
\end{itemize}
\end{multicols}
\end{enumerate}
\end{Exemple}

\newpage



\subsection*{III. Médiatrice d'un segment}
\addcontentsline{toc}{subsection}{III. Médiatrice d'un segment}

\begin{Définition*}[c]
\begin{minipage}[c]{0.65\textwidth}
La \textbf{médiatrice} d'un segment est la droite qui est perpendiculaire à ce segment et qui passe par son milieu.
\end{minipage} \hfill \begin{minipage}[c]{0.35\textwidth}
\begin{center}
\begin{tikzpicture}[scale=0.8]
% Segment AB
\fill[lightgray] (-0.3,0.3) rectangle (0,0);
\draw[gray] (0,0.3) -- (-0.3,0.3);
\draw[gray] (-0.3,0.3) -- (-0.3,0);


\draw[ultra thick,Couleur8] (-2,0) -- (2,0);
\draw[ultra thick, Couleur4] (0,-0.5) -- (0,2);
\node[black] at (-2,0) {\textbf{+}};
\node[black] at (-2,-0.3) {A};
\node[black] at (2,0) {\textbf{+}};
\node[black] at (2,-0.3) {B};

% Médiatrice


\draw[thick,gray] (-1.1,0.15) -- (-1.1,-0.15);
\draw[thick,gray] (-1,0.15) -- (-1,-0.15);
\draw[thick,gray] (1.1,0.15) -- (1.1,-0.15);
\draw[thick,gray] (1,0.15) -- (1,-0.15);
% Label médiatrice
\node[Couleur4, right] at (0,1.8) {Médiatrice};
\end{tikzpicture}
\end{center}
\end{minipage}

\end{Définition*}



\begin{Propriétés*}
\begin{minipage}[c]{0.65\textwidth}
\begin{itemize}[label=\textcolor{Couleur1}{$\bullet$}]
\item Si un point M appartient à la médiatrice $(d)$ d'un segment [AB], alors il est \textbf{équidistant} de ses extrémités. Autrement dit, si M $\in (d)$ alors MA = MB.
\item Si un point n'appartient pas à la médiatrice d'un segment, alors il est plus proche de l'une des extrémités que de l'autre.
\end{itemize}

\end{minipage} \hfill \begin{minipage}[c]{0.35\textwidth}
\begin{center}
\begin{tikzpicture}[scale=0.8]
\draw[dashed, ultra thick,Couleur7] (-2,0) -- (0,1.6);
\draw[dashed, ultra thick,Couleur7] (2,0) -- (0,1.6);
\fill[lightgray] (-0.3,0.3) rectangle (0,0);
\draw[gray] (0,0.3) -- (-0.3,0.3);
\draw[gray] (-0.3,0.3) -- (-0.3,0);


\draw[ultra thick,Couleur8] (-2,0) -- (2,0);
\draw[ultra thick, Couleur4] (0,-0.5) -- (0,2);
\node[black] at (-2,0) {\textbf{+}};
\node[black] at (-2,-0.3) {A};
\node[black] at (2,0) {\textbf{+}};
\node[black] at (2,-0.3) {B};
\node[black] at (0,1.6) {\textbf{+}};
\node[black] at (-0.5,1.8) {M};
\node[Couleur4] at (0.5,1.8) {$(d)$};
% Médiatrice


\draw[thick,gray] (-1.1,0.15) -- (-1.1,-0.15);
\draw[thick,gray] (-1,0.15) -- (-1,-0.15);
\draw[thick,gray] (1.1,0.15) -- (1.1,-0.15);
\draw[thick,gray] (1,0.15) -- (1,-0.15);


\draw[thick,gray] (-1.1,0.95) -- (-0.9,0.65);

\draw[thick,gray] (1.1,0.95) -- (0.9,0.65);


\end{tikzpicture}
\end{center}
\end{minipage}

\end{Propriétés*}


\begin{Exemple} 
\begin{minipage}[c]{0.65\textwidth}
Le point P n'appartient pas à la médiatrice de [AB] : il est plus proche de A que de B.
\end{minipage} \hfill \begin{minipage}[c]{0.35\textwidth}

\begin{center}
\begin{tikzpicture}[scale=0.8]
\draw[dashed, ultra thick,Couleur7] (-2,0) -- (-0.5,1.6);
\draw[dashed, ultra thick,Couleur7] (2,0) -- (-0.5,1.6);
\fill[lightgray] (-0.3,0.3) rectangle (0,0);
\draw[gray] (0,0.3) -- (-0.3,0.3);
\draw[gray] (-0.3,0.3) -- (-0.3,0);


\draw[ultra thick,Couleur8] (-2,0) -- (2,0);
\draw[ultra thick, Couleur4] (0,-0.5) -- (0,2);
\node[black] at (-2,0) {\textbf{+}};
\node[black] at (-2,-0.3) {A};
\node[black] at (2,0) {\textbf{+}};
\node[black] at (2,-0.3) {B};
\node[black] at (-0.5,1.6) {\textbf{+}};
\node[black] at (-1,1.8) {P};
\node[Couleur4] at (0.5,1.8) {$(d)$};
% Médiatrice


\draw[thick,gray] (-1.1,0.15) -- (-1.1,-0.15);
\draw[thick,gray] (-1,0.15) -- (-1,-0.15);
\draw[thick,gray] (1.1,0.15) -- (1.1,-0.15);
\draw[thick,gray] (1,0.15) -- (1,-0.15);




\end{tikzpicture}
\end{center}

\end{minipage}

\end{Exemple}

\begin{Propriété*}
\begin{minipage}[c]{0.65\textwidth}
Si un point K est situé à égale distance de deux points A et B, alors il appartient à la médiatrice $(d)$ du segment [AB]. Autrement dit, si KA = KB alors K $\in (d)$.

\end{minipage} \hfill \begin{minipage}[c]{0.35\textwidth}
\begin{center}
\begin{tikzpicture}[scale=0.8]
\draw[dashed, ultra thick,Couleur7] (-2,0) -- (0,-1.6);
\draw[dashed, ultra thick,Couleur7] (2,0) -- (0,-1.6);
\fill[lightgray] (0,-0.3) rectangle (0.3,0);
\draw[gray] (0.3,0) -- (0.3,-0.3);
\draw[gray] (0.3,-0.3) -- (0,-0.3);


\draw[ultra thick,Couleur8] (-2,0) -- (2,0);
\draw[ultra thick, Couleur4] (0,0.5) -- (0,-2);
\node[black] at (-2,0) {\textbf{+}};
\node[black] at (-2,-0.3) {A};
\node[black] at (2,0) {\textbf{+}};
\node[black] at (2,-0.3) {B};
\node[black] at (0,-1.6) {\textbf{+}};
\node[black] at (-0.5,-1.8) {K};
\node[Couleur4] at (0.5,0.5) {$(d)$};
% Médiatrice


\draw[thick,gray] (-1.1,0.15) -- (-1.1,-0.15);
\draw[thick,gray] (-1,0.15) -- (-1,-0.15);
\draw[thick,gray] (1.1,0.15) -- (1.1,-0.15);
\draw[thick,gray] (1,0.15) -- (1,-0.15);


\draw[thick,gray] (-1.1,-0.95) -- (-0.9,-0.65);

\draw[thick,gray] (1.1,-0.95) -- (0.9,-0.65);


\end{tikzpicture}
\end{center}
\end{minipage}

\end{Propriété*}



\begin{Propriété}\textbf{Caractéristique de la médiatrice}. \\
 
La médiatrice d'un segment est l'ensemble des points équidistants des extrémités de ce segment.
\end{Propriété}

\begin{Exemple}
\vspace{-0.5cm}
\begin{minipage}[c]{0.65\textwidth}
D'après le codage, M, N, P et Q sont à égale distance de A et de B donc ils appartiennent tous à la médiatrice de [AB], tracée ici en rouge.


\end{minipage} \hfill \begin{minipage}[c]{0.35\textwidth}
\begin{center}


\begin{tikzpicture}[scale=1]

% Définition des coordonnées des points
\coordinate (A) at (0,2.5);
\coordinate (B) at (0,-0.5);
\coordinate (Q) at (-0.5,1);
\coordinate (P) at (0.3,1);
\coordinate (N) at (1.2,1);
\coordinate (M) at (2.5,1);

% Lignes depuis A
\draw[very thick, Couleur11] (A) -- (B);
\draw[very thick, Couleur12] (A) -- (Q);
\draw[very thick, Couleur9] (A) -- (P);
\draw[very thick, Couleur7] (A) -- (N);
\draw[very thick, Couleur10] (A) -- (M);

% Lignes depuis B
\draw[very thick, Couleur12] (B) -- (Q);
\draw[very thick, Couleur9] (B) -- (P);
\draw[very thick, Couleur7] (B) -- (N);
\draw[very thick, Couleur10] (B) -- (M);

% Ligne horizontale
\draw[ultra thick, Couleur4] (-1,1) -- (3,1);

% Marques d'égalité des segments

% Pour AQ = QB (un trait perpendiculaire sur chaque segment)
\coordinate (MidAQ) at ($(A)!0.5!(Q)$);
\coordinate (MidQB) at ($(Q)!0.5!(B)$);
\draw[thick, gray] ($(MidAQ)!0.1cm!90:(Q)$) -- ($(MidAQ)!0.1cm!-90:(Q)$);
\draw[thick, gray] ($(MidQB)!0.1cm!90:(B)$) -- ($(MidQB)!0.1cm!-90:(B)$);

% Pour AP = PB (un cercle sur chaque segment)  
\coordinate (MidAP) at ($(A)!0.5!(P)$);
\coordinate (MidPB) at ($(P)!0.5!(B)$);
\draw[thick, gray] (MidAP) circle (0.08cm);
\draw[thick, gray] (MidPB) circle (0.08cm);

% Pour AN = NB (une croix sur chaque segment)
\coordinate (MidAN) at ($(A)!0.5!(N)$);
\coordinate (MidNB) at ($(N)!0.5!(B)$);
\draw[thick, gray] ($(MidAN) + (-0.08,0)$) -- ($(MidAN) + (0.08,0)$);
\draw[thick, gray] ($(MidAN) + (0,-0.08)$) -- ($(MidAN) + (0,0.08)$);
\draw[thick, gray] ($(MidNB) + (-0.08,0)$) -- ($(MidNB) + (0.08,0)$);
\draw[thick, gray] ($(MidNB) + (0,-0.08)$) -- ($(MidNB) + (0,0.08)$);

% Pour AM = MB (double trait perpendiculaire sur chaque segment)
\coordinate (MidAM) at ($(A)!0.5!(M)$);
\coordinate (MidMB) at ($(M)!0.5!(B)$);
% Premier trait pour AM
\draw[thick, gray] ($(MidAM)!0.1cm!90:(M)$) -- ($(MidAM)!0.1cm!-90:(M)$);
% Deuxième trait pour AM (décalé)
\coordinate (MidAM2) at ($(MidAM)!0.12cm!(M)$);
\draw[thick, gray] ($(MidAM2)!0.1cm!90:(M)$) -- ($(MidAM2)!0.1cm!-90:(M)$);
% Premier trait pour MB
\draw[thick, gray] ($(MidMB)!0.1cm!90:(B)$) -- ($(MidMB)!0.1cm!-90:(B)$);
% Deuxième trait pour MB (décalé)
\coordinate (MidMB2) at ($(MidMB)!0.12cm!(B)$);
\draw[thick, gray] ($(MidMB2)!0.1cm!90:(B)$) -- ($(MidMB2)!0.1cm!-90:(B)$);

% Points (croix)
\draw[thick, black] (A) +(-2pt,-2pt) -- +(2pt,2pt);
\draw[thick, black] (A) +(-2pt,2pt) -- +(2pt,-2pt);
\draw[thick, black] (B) +(-2pt,-2pt) -- +(2pt,2pt);
\draw[thick, black] (B) +(-2pt,2pt) -- +(2pt,-2pt);
\draw[thick, black] (Q) +(-2pt,-2pt) -- +(2pt,2pt);
\draw[thick, black] (Q) +(-2pt,2pt) -- +(2pt,-2pt);
\draw[thick, black] (P) +(-2pt,-2pt) -- +(2pt,2pt);
\draw[thick, black] (P) +(-2pt,2pt) -- +(2pt,-2pt);
\draw[thick, black] (N) +(-2pt,-2pt) -- +(2pt,2pt);
\draw[thick, black] (N) +(-2pt,2pt) -- +(2pt,-2pt);
\draw[thick, black] (M) +(-2pt,-2pt) -- +(2pt,2pt);
\draw[thick, black] (M) +(-2pt,2pt) -- +(2pt,-2pt);

% Étiquettes des points
\node[above left] at (A) {A};
\node[below left] at (B) {B};
\node[above left] at (Q) {Q};
\node[below right] at (P) {P};
\node[below right] at (N) {N};
\node[above right] at (M) {M};

\end{tikzpicture}
\end{center}
\end{minipage}

\end{Exemple}

\begin{Remarque}
Il est possible de construire la médiatrice d'un segment uniquement à l'aide d'un compas et d'une règle non graduée : LLS.fr/M6Methode26
\end{Remarque}

